\section*{2 Preliminaries}
We consider multivariate polynomials over a field $\mathbb F$ of characteristic0. A polynomial $p$ is called square-free, if for any non-constant irreducible factor $q$, the polynomial $q^2$ is not a factor of $p$.
By $\deg(p)$ we denote the total degree of $p$. Let $x$ and $\mathbf z=(z_1,\ldots,z_n)$ be variables and $p(x,\mathbf z)$ be a $(n+1)$-variate polynomial. Then we can view $p$ as a univariate polynomial $p=\sum_i a_i(\mathbf z)x^i$ over $K[x]$, where $K=\mathbb F[\mathbf z]$. The $x$-degree of $p$ is denoted by $\deg_x(p)$. It is the highest degree of $x$ in $p$. Polynomial $p$ is called monic in $x$, if the coefficient $a_{d_x}(\mathbf z)$ is a nonzero constant, where $d_x=\deg_x(p)$.
By $\operatorname{poly}(n)$ we denote the class of polynomials in $n\in\mathbb N$. We denote by $I=\langle x\rangle$ the ideal of polynomials generated by $x$ over the ring $\mathbb F[x,\mathbf z]$. The $k$-th power of the ideal $I$ is the ideal $I^k=\langle x^k\rangle$.
\paragraph{Computational models}
An arithmetic circuit is a directed acyclic graph, whose leaf nodes are labeled by the variables $x_1,\ldots,x_n$ and various constants from the underlying field. The other nodes are labeled by sum gates or product gates. A node labeled by a variable or constant computes the same. A node labeled by sum or product compute the sum or product of the polynomials computed by nodes connected by incoming edges. The size of an arithmetic circuit is the total number of its edges.
An arithmetic formula is a special kind of arithmetic circuit. A formula has the structure of a directed acyclic tree. Every node in a formula has out-degree at most one. As we can not reuse computations in a formula, it is considered to be weaker than circuits.
An arithmetic branching program (ABP) is a layered directed acyclic graph with a single source node and a single sink node. An edge of an ABP is labeled by a variable or a constant from the field. The weight corresponding to a path from the source to the sink is the product of the polynomials labeling the edges on the path. The polynomial $f(x_1,\ldots,x_n)$ computed by the ABP is the sum of the weights of the all possible paths from source to sink.
The size of an ABP is the number of its edges. The size of the smallest ABP computing $f$ is denoted by $\sz(f)$. The degree of a polynomial computed by an ABP of size $s$ is at most $\operatorname{poly}(s)$.
\paragraph{Properties of ABPs}
Univariate polynomials have small ABPs. Let $p(x)$ be a univariate polynomial of degree $d$. It can be computed by an ABP of size $2d+1$, actually even by a formula of that size.
For univariate polynomials $p(x),q(x)$, the extended Euclidian algorithm computes the gcd $h=\gcd(p,q)$ and also the B\'ezout-coefficents, polynomials $a,b$ such that $ap+bq=h$, where $\deg(a)<\deg(q)$ and $\deg(b)<\deg(p)$. Let $p$ have the larger degree, $d=\deg(p)\geq\deg(q)$. Then clearly also $\deg(h),\deg(a),\deg(b)\leq d$, and consequently, all these polynomials,$p,q,h,a,b$ have ABP-size at most $2d+1$.
Let $p(\mathbf x),q(\mathbf x)$ be multivariate polynomials in $\mathbf x=(x_1,\ldots,x_n)$. For the ABP-size with respect to addition and multiplication, we have
\begin{enumerate}\item $\sz(p+q)\leq\sz(p)+\sz(q)$,\item $\sz(pq)\leq\sz(p)+\sz(q)$.\end{enumerate}
For the sum of two ABPs $B_p,B_q$ one can put $B_p$ and $B_q$ in parallel by merging the two source nodes of $B_p$ and $B_q$ into one new source node, and similar for the two sink nodes. For the product, one can put $B_p$ and $B_q$ in series by merging the sink of $B_p$ with the source of $B_q$.
Another operation is substitution. Let $p(x_1,\ldots,x_n)$ and $q_1(\mathbf x),\ldots,q_n(\mathbf x)$ be polynomials. Let $\sz(q_i)\leq s$, for $i=1,\ldots,n$. Then we have
\[\sz(p(q_1,\ldots,q_n))\leq s\cdot\sz(p).\]
To get an ABP for $p(q_1(\mathbf x),\ldots,q_n(\mathbf x))$, replace an edge labeled $x_i$ in the ABP $B_p$ for $p$ by the ABP $B_{q_i}$ for $q_i$.
It is known that the determinant of a symbolic matrix of dimension $n$ can be computed by an ABP of size $\operatorname{poly}(n)$~\cite{ref15}. By substitution, the entries of the matrix can itself be polynomials computed by ABPs.
\paragraph{Resultant}
Given two polynomials $p(x,\mathbf y)$ and $q(x,\mathbf y)$ in variables $x$ and $\mathbf y=(y_1,\ldots,y_n)$, consider them as polynomials in $x$ with coefficients in $\mathbb F[\mathbf y]$. The resultant of $p$ and $q$ w.r.t.$x$, denoted by $\Res_x(p,q)$, is the determinant of the Sylvester matrix of $p$ and $q$. For the definition of the Sylvester matrix, see~\cite{ref20}. Note that $\Res_x(p,q)$ is a polynomial in $\mathbb F[\mathbf y]$.
Basic properties of the resultant are that it can be represented as a linear combination of $p$ and $q$, and that it provides information about the gcd of $p$ and $q$.
\setcounter{nttheorem}{1}
\begin{ntlemma}[Author Lemma 2, See {\cite{ref20}}]
Let $p(x,\mathbf y)$ and $q(x,\mathbf y)$ be polynomials of degree $\leq d$ and $h=\gcd(p,q)$.
\begin{enumerate}\item $\deg(\Res_x(p,q))\leq4d^2$,\item $\exists u,v\in\mathbb F[x,\mathbf y]\ up+vq=\Res_x(p,q)$,\item $\Res_x(p,q)=0\Longleftrightarrow\deg_x(h)>0$.\end{enumerate}
\end{ntlemma}
Note that the problem whether $\Res_x(p,q)=0$ is a polynomial identity test (PIT), because $\Res_x(p,q)\in\mathbb F[\mathbf y]$. It can be solved in a randomized way by the DeMillo-Lipton-Schwartz-Zippel Theorem (see~\cite{ref4} and the references therein for more details and history of this theorem).
\begin{nttheorem}[Author Theorem 3: Polynomial Identity Test]
Let $p(\mathbf x)$ be an $n$-variate nonzero polynomial of total degree $d$. Let $S\subseteq\mathbb F$ be a finite set. For $\boldsymbol\alpha\in S^n$ picked independently and uniformly at random,
\[\Pr[p(\boldsymbol\alpha)=0]\leq\frac d{|S|}.\]
\end{nttheorem}
